\section{近期研究（2026 年 8--11 月）}
\paragraph{CLEAR：让低分辨率观测和超分先验各自约束可信区域}
2026.08~\cite{2608_02206}。少视图超分常先用二维模型为每张图生成更清晰的参考，但补出的细节可能与真实低分辨率观测或其他视图冲突。CLEAR 在同一个高斯场内共同优化这两类信号，按高斯处理相冲突的梯度。它把图像局部的可靠超分证据映射到可见高斯，只让证据较强的高频信号驱动局部更新和增密。

原文在三种四倍超分基准上评估，包括 LLFF 三输入视图和 Blender 八输入视图。真实低分辨率图和由冻结二维模型生成的高分辨率参考对应不同监督来源，不能将生成的局部细节视作实测结构。结果也不直接覆盖不同倍率或未见过的退化方式。

\paragraph{JSGS：按每帧 JPEG 量化表匹配渲染观测}
2026.08~\cite{2608_08659}。高斯场若直接拟合经过不同强度 JPEG 压缩的训练照片，容易把块状伪影变成共享的场景结构。JSGS 读取每张训练图的量化表。图像域中，它先按对应 JPEG 状态处理渲染，再与已解码的输入比对；频域中，它按亮度量化表调整离散余弦变换残差的权重，再将不一致的证据用于原语正则。

评测在四个 LLFF 场景和三个 Mip-NeRF 360 场景进行，仅压缩训练视图，并保留无压缩的测试视图。三种混合质量设置下，JSGS 的平均 SSIM 和 LPIPS 最优，但 FDS-GS 保持较高平均 PSNR。这些指标体现不同取舍。方法也未处理双重 JPEG，因为文件头仅记录最后一轮的量化状态。

\paragraph{3DGART：按高斯原语收集光追反向梯度}
2026.08~\cite{2608_17298}。在可微高斯光追中，许多像素的反向计算可能同时写入同一个高斯的梯度，造成写入争用。3DGART 将反向过程按瓦片和高斯组织，经由线程在所属瓦片内汇总原语对多个像素的梯度贡献，再更新原语，减少大量像素线程竞争同一参数。这里优化的是反传工作分配，正向仍沿射线查询高斯。

论文在 Mip-NeRF 360 上报告相对 3DGRT 约四倍的训练加速，并在标准新视图数据上比较重建。该实现用部分显存开销换取梯度汇总速度，混合反传用于缓解内存成本。使用的透视投影也限制了对鱼眼等非线性相机的直接支持。

\paragraph{FixAnything：从渲染序列而非单帧修复视图}
2026.08~\cite{2608_23549}。FixAnything 把粗糙三维场景沿已知相机路径渲染成视频，再用视频扩散补全伪影。输入掩码标记可信参考帧，使修复器保留场景线索，同时利用序列条件维持相机运动。训练还通过偏好优化强化与目标相机轨迹一致的修复结果。方法作用于渲染视频，因此输入表示可以是高斯、NeRF、网格或稀疏点云。

同一个模型在 DL3DV 的 3、6、9 个输入视角设置下对照四类初始表示。报告的效果是修复渲染质量，不等于修改了输入表示的几何。模型在推理时按重叠的 61 帧片段执行 50 步去噪，部署成本与直接高斯渲染不同。

\paragraph{WilLaGS：把图像外观变化与临时遮挡分开建模}
2026.08~\cite{2608_28240}。同一地点的照片可能跨越不同时间和天气。WilLaGS 用 β-VAE 学出连续外观潜空间，并将每幅图的潜变量送入三维神经外观场，赋予光照与外观空间变化能力。另一条教师学生通路根据感知差异估计不一致区域，降低移动物体等临时成分对场景几何的影响。

论文在 Photo Tourism 的地标及 NeRF-OSR 的户外光照场景上比较重建。Photo Tourism 的 Sacre Coeur 场景报告 25.84 dB，比同表中 AsymGS 高 2.28 dB。这是一个场景的对照，不能替代跨场景的完整指标；该掩码的作用是减弱与静态解释不一致的像素，并未为这些移动物体建立可编辑三维模型。

\paragraph{TruncGradGS：扩大光栅化反传的有效梯度范围}
2026.09~\cite{2609_03534}。光栅化时，离原语中心较远的像素给出的高斯导数接近零，因此未初始化到表面附近的原语很难移动到正确位置。TruncGradGS 在反向过程中用连续的分段代理替代近零的尾部梯度，使更远的像素也参与更新。前向高斯表示及推理渲染不因这个训练策略改变。

在 Mip-NeRF 360 和新数据集的静态子集，论文分别用随机与 COLMAP 初始化评估 3DGS 和 2DGS，并另评估动态场景。收益在随机初始化下尤其明显。反向加入更多像素与原语配对通常会增加训练时间，因此这种改动改善的是优化可达性，不是无额外训练成本的渲染加速。

\paragraph{TileGS：把全局深度排序替换为局部深度分桶}
2026.09~\cite{2609_03613}。TileGS 在每个瓦片内按深度分桶，再对需要的局部情况作精确修复。它保留高斯投影输入和合成语义，改变的是栅格化核读取排序任务的粒度。

在九个场景和桌面及笔记本 Ada GPU 上，作者固定投影输入与 gsplat 对照，并分开报告整管线和核函数开销。主要结果用 No-GW 与 64 桶配置，硬件及桶数会影响最终收益。

\paragraph{RouteBridge：让 NeRF 和 3DGS 按射线互相蒸馏}
2026.09~\cite{2609_09606}。NeRF 和高斯场在不同区域可能各有优势。RouteBridge 为两种表示保留各自的训练和部署接口，只通过射线颜色、不透明度与深度进行蒸馏。路由器逐射线比较光度和几何线索，决定由哪一侧提供监督，证据不足时允许弃权。

论文在 Mip-NeRF 360 将原始 3DGS 对照的 27.21 dB 提高到导出高斯的 28.77 dB，并对静态三视图 DTU 报告 21.12 dB。结果是所测训练配置下的单模型导出质量。训练仍需同时准备两种表示，若两侧的几何估计都不可靠，可靠度路由也不能补出实测信息。

\paragraph{用学习到的历史置信度去除随机高斯渲染噪声}
2026.09~\cite{2609_25604}。随机高斯泼溅以采样代替确定性排序，降低部分渲染成本的同时可能产生噪声与闪烁。该研究在两路指数移动平均历史上估计当前像素可复用的历史信息，以各向异性邻域滤波和基于方差的选择输出较稳的结果。历史复用过强会拖出重影，复用不足则保留闪烁。

原文在 MipNeRF360 主发布的七个场景和 Tanks and Temples 子集测量静态及相机轨迹表现，图像宽度设为 1600 像素，训练原语上限为 300 万。作者说明最终仍可能闪烁，稳定项以少量重影换取抑制。该实验范围不能推出运动物体上的时间一致性。

\paragraph{用主眼重投影与局部补渲染加速双目高斯显示}
2026.09~\cite{2609_30741}。双目显示常对左右眼分别执行完整光栅化，但两眼只存在有限视差。该方法完整渲染主眼后，用加权深度代理将颜色重投影到副眼。小孔洞由插值填补，大孔洞才调用局部高斯光栅化，从而减少第二次全帧渲染。两种比较路径均使用相同的固定 2DGS 场景与主眼渲染。

在 DTU、Tanks and Temples 和 Mip-NeRF 360 的静态双目对照中，完整双目渲染时间下降 15.5\% 至 28.8\%，显存下降 6\% 至 11\%，副眼 PSNR 则损失 0.82 至 1.20 dB。测试只在副眼有拍摄参考，不包含真实头显主观阈值或连续头动，因此不能将画质损失判断为人眼不可察觉。

\paragraph{ControlGS：让高斯表示适应渲染后的显示管线}
2026.09~\cite{2609_32038}。头显最终显示的图像还会经过抗锯齿重采样、镜头校正和显示映射，直接优化渲染器输出未必能兼顾这些变化。ControlGS 将三类后处理状态作为条件输入，在神经高斯解码时调节生成的原语，并以模拟显示光学后的画面作为端到端目标。

原文用 Mip-NeRF 360 与 NeRF Synthetic 比较多种神经高斯骨干，汇总 20 种运行设置，并提供对照各方法都使用超采样的结果。该协议模拟显示光路，未覆盖大部分人视觉系统特性或真实头显主观质量；换设备或更改处理链时可能需要重新训练或校准。

\paragraph{Gaussian Stippling：将随机可见性采样和时空去噪合起来}
2026.09~\cite{2609_38488}。该方法在固定的训练后高斯资产上混合按原语和按片元的随机采样，使用当前和历史证据重建低采样渲染结果，从而避免传统排序路径。改变的是可见性取样与图像恢复过程。

论文用三个基准的 11 个场景评价，默认针对各场景用训练图像训练重建网络，再在该场景的保留视图上测量。因而应把按场景训练的性能和直接跨场景部署的性能分开比较。

\paragraph{EffGS：按频段和投影范围分配高斯更新}
2026.09~\cite{2609_39553}。统一的图像误差和全局增密会使原语投入与细节需求不匹配。EffGS 将重建残差与分阶段的高斯差分掩码组合，按频率调整像素权重；增密和剪枝只作用于被采样视图采到且具有有效投影范围的原语。各原语还有可学习的尺度调节量，用来适配局部有效支持。

原文在 Mip-NeRF 360、Deep Blending 和 Tanks and Temples 及城市级数据比较画质、训练时间和参数量，另提供匹配原语预算实验与沿用 CityGS 切块策略的多 GPU 扩展。相同原语数并不代表相同算力或延迟，实验主要覆盖静态场景。

\paragraph{TileSkipper：离线编译逐高斯的瓦片裁剪阈值}
2026.10~\cite{2610_09343}。TileSkipper 接收已训练高斯检查点，把光栅器的全场景贡献阈值换成逐原语的静态策略。离线阶段按候选阈值的瓦片对节省和近似失真作分配，运行时只读取每个高斯的一字节标记及查找表，不更新检查点参数也不逐帧预测阈值。

原文在 13 个场景和十种模型后端集成比较。对已依赖不透明度边界的六条渲染路径，所报额外加速约为 1.009 至 1.121 倍；与 AdaGScale 的原始分辨率比较区间跨越一倍，不能说明普遍更快。编译需要训练相机和图像来估计节省与误差，因此不能直接用于任意只公开了点云的检查点。

\paragraph{在高斯预算内按图像置信度增密和重定位}
2026.10~\cite{2610_13541}。3DGS-MCMC 用固定预算逐步增长原语，但按不透明度选址与回收未必对应当前渲染误差。该方法改用置信度加权的图像误差决定从哪里增加原语，再依它对可见图像的贡献决定哪些已有高斯可以重定位，保留预算和渐进增长流程。

论文在多个静态基准和 10 万到 200 万原语预算上对照，所测预算相对 3DGS-MCMC 均获改善，但最低预算仍明显取决于初始化。由 SfM 点启动时，10 万高斯的结果仍落后于 ConeGS；即使采用相同初始化，ConeGS 的 LPIPS 仍较好。因此增密策略不能取代紧预算下的可靠初值。

\paragraph{AGSR：用不同投影孔径的高斯渲染指导超分}
2026.11~\cite{2611_25934}。在低分辨率训练的高斯场上直接提高输出分辨率，不能保证恢复高频细节。AGSR 保持高斯表示固定，缩小协方差得到不同投影孔径的 RGB、透射率和深度。跨孔径差异作为结构不确定性提示，逐级网络融合多组观测后进行重建。

原文在 Mip-NeRF 360 比较两倍与四倍超分，以及不同训练图降采样设置。模型可按需要使用一至三种孔径；输出分辨率和退化程度不同，训练和推理条件也不同。实验限于该基准，不能用最轻配置声称所有超分组合都有相同延迟。
